\chapter{Coordinates and Vectors}
\label{ch:ch02-coordinates-vectors}

\section{Why Position Is Not Enough}

The previous chapter studied similarity and scale, and found that certain
relationships survive changes in size: a triangle may grow or shrink while
its shape, recorded in the ratios between its sides, remains unchanged.
Geometry must answer a different question as well. A hiker begins at one
location and ends at another; a ship leaves a harbor and travels across open
water; a package is moved across a warehouse floor. In each case what
matters is not merely where something is, but how it moved. A single point
describes a location, but it does not describe a displacement, and to
describe movement requires a mathematical object capable of representing
magnitude and direction at once. The development of such an object begins,
as it must, with a way of describing position in the first place.

\section{Coordinates as a System of Description}

\begin{definition}
The \emph{Cartesian coordinate system} consists of two perpendicular number
lines meeting at a point called the \emph{origin}: a horizontal
\emph{$x$-axis} and a vertical \emph{$y$-axis}. Every point in the plane is
assigned an ordered pair $(x,y)$, where $x$ measures horizontal displacement
from the origin and $y$ measures vertical displacement.
\end{definition}

\begin{historicalnote}{Why ``Cartesian''}
The coordinate system of this chapter is named for the French philosopher
and mathematician Ren\'e Descartes, whose Latinized name,
\emph{Cartesius}, gives the adjective \emph{Cartesian} \cite{boyer1991}.
Descartes did not present the system in the clean $(x,y)$-axis form given
here; his 1637 appendix \emph{La G\'eom\'etrie} used a single reference
line and measured a second distance obliquely from it, arriving at
essentially the same idea, that a curve could be described by an
algebraic relationship between two measured quantities, without drawing
the perpendicular axes now named after him. The now-standard picture of
two perpendicular axes meeting at an origin was assembled gradually by
later mathematicians working from Descartes' insight, so that the name
Cartesian honors the founding idea, the translation of geometric
questions into algebraic ones, more than it records the exact picture
students draw today.
\end{historicalnote}

The pair $(3,2)$, for instance, identifies the point three units to the
right of the origin and two units above it. Coordinates should not be
mistaken for part of the plane itself; they are a convention for labeling
positions within it, and different coordinate systems may describe the same
geometric reality in different ways. The usefulness of Cartesian coordinates
lies in the fact that geometric questions can be translated directly into
algebraic ones, a translation this chapter now puts to use.

\section{Vectors}

Consider the points $A = (1,1)$ and $B = (4,3)$. Moving from $A$ to $B$
requires a horizontal change of $4 - 1 = 3$ units and a vertical change of
$3 - 1 = 2$ units, and this displacement is recorded by the ordered pair
$\langle 3, 2 \rangle$.

\begin{definition}
A \emph{vector} is an object possessing magnitude and direction but no fixed
location. Two vectors are \emph{equal} whenever they share the same
magnitude and direction, regardless of where each is drawn.
\end{definition}

The displacement $\langle 3,2 \rangle$ describes the same movement whether
it begins at $A$, at $B$, or at any other point of the plane, which is
precisely why a vector is not tied to a particular location the way a point
is. For this reason vectors are often drawn as arrows, with the length of
the arrow representing magnitude and its orientation representing direction.

\section{Vector Addition}

Suppose one displacement is $\mathbf{u} = \langle 2,1 \rangle$ and another is
$\mathbf{v} = \langle 3,2 \rangle$. Performing the first displacement and
then the second produces the combined displacement
\[
\mathbf{u} + \mathbf{v} = \langle 2+3,\ 1+2 \rangle = \langle 5,3 \rangle.
\]
Geometrically, this sum may be visualized by placing the tail of the second
vector at the head of the first, so that the resulting displacement runs
from the tail of $\mathbf{u}$ to the head of $\mathbf{v}$; this construction
is known as the \emph{tip-to-tail method}, and an equivalent picture is
provided by the parallelogram rule.

\section{Scalar Multiplication}

Chapter~1 studied scale transformations of the form $(x,y) \mapsto (kx,ky)$,
and vectors behave in exactly the same way under scaling. Given a vector
$\mathbf{v} = \langle x,y \rangle$ and a scalar $k$, the product
$k\mathbf{v} = \langle kx, ky \rangle$ points in the same direction as
$\mathbf{v}$ when $k>0$, with $3\langle 2,1\rangle = \langle 6,3\rangle$
illustrating a vector three times as long as the original. If $0 < k < 1$
the vector becomes shorter without changing direction, and if $k < 0$ the
vector reverses direction entirely. Scalar multiplication is therefore
nothing more than the dilation of Chapter~1, applied to a displacement
rather than to a whole figure.

\section{Derivation: The Length of a Vector}

A vector possesses both direction and magnitude, and this section derives a
formula for the latter. Consider $\mathbf{v} = \langle x,y \rangle$ with its
tail placed at the origin. The horizontal and vertical components of
$\mathbf{v}$, together with $\mathbf{v}$ itself, form a right triangle whose
legs have lengths $|x|$ and $|y|$ and whose hypotenuse is $\mathbf{v}$.

\begin{theorem}
For a vector $\mathbf{v} = \langle x,y \rangle$, its magnitude is
\[
|\mathbf{v}| = \sqrt{x^2+y^2}.
\]
\end{theorem}

\begin{proof}
The vector $\mathbf{v}$ forms the hypotenuse of a right triangle whose legs
have lengths $|x|$ and $|y|$. By the Pythagorean theorem,
$|\mathbf{v}|^2 = x^2 + y^2$, and taking the positive square root gives
$|\mathbf{v}| = \sqrt{x^2+y^2}$.
\end{proof}

This formula is one of the most important results of the chapter: it will
later become the bridge connecting vectors to the trigonometric ratios of
Chapter~3, since every vector, by virtue of possessing a horizontal and a
vertical component, already carries a right triangle within it.

\section{Components and Decomposition}

Every vector in the plane may be viewed as the sum of a horizontal
displacement and a vertical displacement. For instance,
\[
\langle 5,3 \rangle = \langle 5,0 \rangle + \langle 0,3 \rangle,
\]
and these two simpler vectors are called the \emph{components} of the
original. The process of expressing a vector as a sum of its components is
called \emph{decomposition}, and it allows complicated motion to be analyzed
one direction at a time; many problems in physics, engineering, and
navigation become tractable only after such a decomposition has been made.

\begin{algorithmproc}{How to Decompose a Vector into Components}
Given a vector $\mathbf{v} = \langle x,y \rangle$, identify the horizontal
coordinate $x$ and the vertical coordinate $y$ separately. Form the
horizontal component $\langle x,0 \rangle$ and the vertical component
$\langle 0,y \rangle$, and confirm that their sum recovers the original
vector, $\mathbf{v} = \langle x,0 \rangle + \langle 0,y \rangle$.
\end{algorithmproc}

\begin{example}
An aircraft travels $120$ kilometers east and $50$ kilometers north. Its
displacement is the vector $\mathbf{v} = \langle 120,50 \rangle$, and by the
magnitude formula,
\[
|\mathbf{v}| = \sqrt{120^2+50^2} = \sqrt{14400+2500} = \sqrt{16900} = 130.
\]
The aircraft is displaced $130$ kilometers from its starting point.
\end{example}

\begin{practicenow}
A hiker travels $24$ kilometers east and $7$ kilometers north. Represent the
displacement as a vector and determine its magnitude.
\end{practicenow}

\section{Applications}

Vectors provide a mathematical language for describing motion, force,
velocity, acceleration, and displacement. Navigation describes courses and
bearings with vectors; physics represents the forces acting on an object
with them; engineering describes stresses and motions within structures in
the same terms; computer graphics determines the positions and movements of
objects within virtual environments the same way still. A vector's
magnitude can already be computed from its components, but its orientation
remains, for now, unspecified: no formal method for measuring direction has
yet been introduced, even though the need for one should already be
apparent. Supplying that method is the task of the next chapter, and the
tool it introduces is the angle.

\section{Looking Ahead}

This chapter introduced vectors as mathematical descriptions of
displacement, and showed how to add them, scale them, decompose them into
components, and determine their magnitudes. Most importantly, it showed that
every vector naturally generates a right triangle, an observation that is
not accidental: right triangles are the bridge between length and
direction. The next chapter begins the study of angles and the ratios
associated with right triangles, ratios that will eventually describe a
vector's direction as precisely as its magnitude can already be described
today. The transition from geometry to trigonometry begins with a single
question: if a vector's magnitude can be measured, how can its direction be
measured?

\section{Exercises}
\begin{exercise}
Find the magnitude of each vector: (a) $\langle 3,4 \rangle$, (b)
$\langle 5,12 \rangle$, (c) $\langle 8,15 \rangle$.
\end{exercise}

\begin{exercise}
Decompose each vector into horizontal and vertical components: (a)
$\langle 7,2 \rangle$, (b) $\langle -4,5 \rangle$, (c) $\langle 3,-6
\rangle$.
\end{exercise}

\begin{exercise}
A boat travels $15$ kilometers east and $20$ kilometers north. Represent the
displacement as a vector and determine its magnitude.
\end{exercise}

\begin{exercise}
Explain why two vectors with the same magnitude and direction are considered
equal even when they begin at different points.
\end{exercise}

\begin{exercise}
A vector is multiplied by a scalar $k$. Explain how its magnitude and
direction change as $k$ varies, and describe what happens in the particular
case $k < 0$.
\end{exercise}
